%%%%%%%%%%%%%%%%%%%%%%%%%%%%%%%%%
% Preliminaries
%%%%%%%%%%%%%%%%%%%%%%%%%%%%%%%%%

\section{Preliminaries}
\label{sec:prelim}

We briefly restate the elements of the Bernoulli set model~\cite{bernoulliSets} needed in this paper.

\begin{axiom}[Element-wise independence]
\label{asm:element_indep}
For all distinct $x, y \in U$, the error events at $x$ and $y$ are mutually independent:
\begin{equation}
P(\mathbf{1}_{\tilde{A}}(x) \neq \mathbf{1}_{A}(x) \mid \mathbf{1}_{\tilde{A}}(y) \neq \mathbf{1}_{A}(y)) =
	P(\mathbf{1}_{\tilde{A}}(x) \neq \mathbf{1}_{A}(x)).
\end{equation}
More generally, the family $\{\mathbf{1}_{\tilde{A}}(x) : x \in U\}$ is mutually independent given the objective set $A$ and the error rates.
Furthermore, within each partition block $B_i$, the error indicators $\{\mathbf{1}_{\tilde{A}}(x) : x \in B_i\}$ are identically distributed.
\end{axiom}

\begin{proposition}
\label{prop:element_rates}
By \cref{asm:element_indep}, for any single element $x \notin A$,
$\Prob{\mathbf{1}_{\tilde{A}}(x) = 1} = \fprate$, and for $x \in A$,
$\Prob{\mathbf{1}_{\tilde{A}}(x) = 1} = \tprate$.
The expected sample rates satisfy $\Expect{\alpha} = \fprate$ and $\Expect{\beta} = \tprate$.
\end{proposition}
\begin{proof}
See~\cite{bernoulliSets}.
\end{proof}

\begin{theorem}
\label{thm:fpbinom}
The random number of false positives $\FP$ given $\RV{N} = n$ in the second-order random approximate set $\tilde{\RV{R}}[\fprate]$ is given by
\begin{equation}
	\FP_n = n \alpha
\end{equation}
with a distribution given by
\begin{equation}
    \FP_n \sim \bindist(n, \fprate).
\end{equation}
\end{theorem}
\begin{proof}
See~\cite{bernoulliSets}.
\end{proof}

\begin{theorem}
\label{thm:fpr}
The random false positive rate $\alpha$ conditioned on $\RV{N} = n$ is denoted by $\alpha_n$ and has a distribution given by
\begin{equation}
    \alpha_n = \frac{\FP_n}{n},
\end{equation}
with an expectation $\fprate$, variance $\fprate(1-\fprate) / n$, and probability mass function
\begin{equation}
	p_{\alpha_n}(\fprateob | \fprate) = p_{\FP_n}(\fprateob n | \fprate)
\end{equation}
over the support $\SetBuilder{ \frac{j}{n} \in \RatSet}{j \in \{0,\ldots,n\}}$.
\end{theorem}
\begin{proof}
See~\cite{bernoulliSets}.
\end{proof}

\begin{corollary}
\label{cor:tnbinom}
Given $n$ negatives, the number of \emph{true negatives} in a random approximate set with a false positive rate $\fprate$ is a random variable denoted by $\TN_n$ with a distribution given by
\begin{equation}
	\TN_n = n - \FP_n \sim \bindist\!\left(n, 1-\fprate\right).
\end{equation}
By definition, the \emph{true negative rate} $\mathrm{TNR}_n = \TN_n / n = 1 - \alpha_n$.
\end{corollary}
\begin{proof}
See~\cite{bernoulliSets}.
\end{proof}

\begin{theorem}
\label{thm:fnbinom}
Given $p$ positives, the uncertain number of \emph{false negatives} in random approximate sets with a false negative rate $\fnrate$ is modeled as a binomial distributed random variable denoted by $\FN_p$,
\begin{equation}
	\FN_p \sim \bindist(p, \fnrate).
\end{equation}
\end{theorem}
\begin{proof}
See~\cite{bernoulliSets}.
\end{proof}

\begin{corollary}
\label{cor:tpbinom}
Given $p$ positives, the number of \emph{true positives} in an approximate set
with a false negative rate $\fnrate$ is a random variable denoted by $\TP_p$
with a distribution given by
\begin{equation}
    \TP_p \sim \bindist(p, \tprate).
\end{equation}
By definition, the \emph{true positive rate} is given by $\TPR_p = 1 -
\beta_p$.
\end{corollary}
\begin{proof}
See~\cite{bernoulliSets}.
\end{proof}
